\section{Instance-Level Semantics Across Modalities}
\label{sec:instance}
Building on the query formulation, this section examines how to realize instance-level semantics. In particular, it addresses the two key challenges identified in Section~\ref{sec:intro}: supporting counterfactual queries efficiently at scale and enabling principled reasoning across heterogeneous data modalities. 

The counterfactual query system must retrieve instances that satisfy feasibility constraints while remaining competitive under the cost induced by $\omega$. This depends strongly on modality: admissibility may be checked via explicit attributes or structural rules, or may require execution- or semantics-level validation. We next analyze how P1 (interpretability), P2 (cost faithfulness), and P3 (efficiency) guide design across tabular, temporal, graph, text, and image data, illustrating patterns that generalize to other data types.

\paragraph{Tabular and Relational Data.}
Among all modalities, tabular data has the most explicit instance-level structure. Since features are represented column-wise, edits have clear coordinate semantics: $\Delta(\mathbf{x}_q,\mathbf{x})$ can be directly obtained via column-wise comparison, the cost $\omega(\mathbf{x}_q,\mathbf{x})$ can aggregate per-column modification effort, and immutable attributes can be enforced by assigning infinite cost. This makes tabular domains the clearest setting for P1 and P2. 

However, counterfactual retrieval is still not equivalent to nearest-neighbor search. Different columns have different semantic scales and actionability, so naive geometric metrics can misalign with true intervention cost~\cite{karabulut2019weighted}. Moreover, immutable attributes and structural rules partition the space into constrained subspaces~\cite{galhotra2022hyper}, and geometric distance serves only as a heuristic rather than a guarantee of feasibility~\cite{bartal2011fast}. Effective systems must therefore encode column-level semantics into indexing and candidate generation, focusing search on viable regions instead of repeatedly exploring infeasible candidates.

\paragraph{Sequence and Temporal Data.}
In temporal or sequential data, the query instance $\mathbf{x}_q$ is a trajectory $(\mathbf{x}_{q,1},\dots,\mathbf{x}_{q,T})$. Transformation cost is often estimated using alignment-based measures such as dynamic time warping (DTW)~\cite{senin2008dynamic} or learned sequence embeddings~\cite{nalmpantis2019signal2vec}, which capture similarity in observation space. However, admissibility is determined by the underlying dynamics rather than by alignment alone. For example, in chess, two move sequences may appear similar yet lead to invalid states because each move depends on the evolving board configuration. 

This separation between similarity and feasibility introduces key challenges. First, edits must correspond to valid trajectory-level transformations, making traceable edits inherently execution-aware~\cite{jing2024towards}. Second, alignment-based similarity is sensitive to time shifts or scaling, leading to poor estimates of true intervention cost~\cite{rajput2024evaluating}. Third, early changes propagate through the sequence, making feasibility path-dependent and requiring execution-level validation~\cite{page2006path}. As a result, systems must go beyond sequence similarity to incorporate dynamics-aware constraints and verification.

\paragraph{Text, Image, and Graph Data.}
Text and image data exhibit a clear separation between scalable retrieval and semantic admissibility. In text, surface metrics such as edit distance or BLEU~\cite{papineni2002bleu} fail to capture semantic change, so the cost induced by $\omega$ is typically approximated in pretrained representation spaces~\cite{devlin2019bert}. However, admissibility depends on higher-level semantics, and local edits may affect global meaning~\cite{yin2015multigrancnn}. As a result, representation proximity can misalign with validity~\cite{wang2021robustness}, and feasibility often requires full-sequence reasoning or external grounding~\cite{gu2024survey}. 

In images, retrieval also relies on embeddings, but pixel-level changes rarely correspond to meaningful semantic edits~\cite{narodytska2017simple}. This weakens traceability~\cite{assran2023self}, misaligns perceptual similarity with intervention cost~\cite{zhang2021multi}, and often requires localization or scene-level reasoning to verify admissibility~\cite{gagliardi2025region}. Recent approaches introduce intermediate semantic layers, such as concept bottleneck models, to make costs and constraints more interpretable~\cite{de2025causally,koh2020concept,dominici2025counterfactual}.

In graph data, instances are structured objects with nodes and edges, where counterfactuals correspond to structural edits~\cite{gao2010survey}, whose comparison under graph edit distance is costly. Retrieval therefore relies on embeddings for scalability~\cite{kipf2016semi}, but these do not capture structural validity, which is governed by constraints such as connectivity, typing, or capacity. As in temporal data, similarity is geometric while validity is structural~\cite{kazemi2019time2vec}, leading to embedding--constraint misalignment~\cite{armgaan2024graphtrail} and costly verification~\cite{xu2018powerful}, since even local edits may induce global changes. Effective graph counterfactuals thus require integrating structural constraints into retrieval.

\paragraph{Extended to Multi-Modality.}
When instances span multiple modalities, counterfactual querying requires coordinated reasoning. A candidate $\mathbf{x}$ must satisfy modality-specific constraints while remaining consistent across modalities~\cite{chen2019cross}. This setting amplifies the core principles: traceability (P1) links edits across representations, cost faithfulness (P2) aligns intervention effort across views, and efficiency (P3) must scale across heterogeneous structures. Moreover, a single semantic factor may manifest in different forms across modalities, complicating admissibility and cost modeling.

In this context, similarity search in joint or modality-specific embeddings serves only as a proposal mechanism, while correctness is determined through cross-modal validation that enforces feasibility and consistency. As a result, multi-modal counterfactual querying is inherently verification-driven. These principles extend beyond single modalities, providing a unified foundation for scalable and consistent counterfactual reasoning across heterogeneous data.

\rev{The constraint type, edit representation, and enforcement point for each modality are summarized in Appendix~\ref{app:eval} (Table~\ref{tab:modality}). For multimodal queries, $\omega$ additionally incorporates a cross-modal consistency predicate: a candidate must satisfy each modality's constraints \emph{and} be jointly consistent across modalities~\cite{chen2019cross} (e.g., a tabular attribute edit must imply a consistent change in the associated text and image views). The current prototype realizes these constraints through validation after shared-embedding retrieval. Crucially, data-centric counterfactual analysis exposes them as \emph{declarative query semantics}: the abstraction enables future optimizers to push modality-specific and cross-modal constraints into indexing and retrieval, analogous to how relational databases evolved from filtering after scans to predicate-aware query optimization. Achieving this---rather than post-filtering---is a central open challenge of this vision.}
